% Textuais/Capitulo10_TeoriaGrafos.tex
\chapter{TEORIA DOS GRAFOS \& REDES COMPLEXAS NA GESTÃO DE FROTAS PÚBLICAS}
\label{cap:grafos}

A modelagem de redes articulou a rede bipartida entre 7.291 veículos e 1.039 oficinas credenciadas, detectando \textbf{18 comunidades regionais com modularidade Louvain $Q = 0,5646$}, além do grafo de co-ocorrência de falhas com $Q = 0,384$.

\begin{figure}[H]
    \centering
    \caption{Rede Bipartida de Estabelecimentos Credenciados e Comunidades de Louvain}
    \includegraphics[width=0.85\textwidth]{figuras/fig94_rede_fornecedores_louvain.png}
    \vspace{0.2cm}
    \fonte{Dados empíricos da frota pública de Santa Catarina processados pelo autor (2026).}
    \label{fig:grafos}
\end{figure}

Ressalva-se que a identificação de agrupamentos densos opera estritamente como ferramenta de priorização amostral para auditoria, e não como comprovação de conluio.
